% ==STYLE-BEGIN==
\documentclass[UTF8,11pt]{ctexart}
\usepackage{amsmath,amssymb,amsthm}
\usepackage[a4paper,margin=2.4cm,top=2.3cm,bottom=2.3cm]{geometry}
\usepackage{fancyhdr}
\usepackage{enumitem}
\usepackage{hyperref}
\hypersetup{hidelinks}

% ---- 版式 ----
\linespread{1.15}
\setlength{\parskip}{0.35em}
\setlength{\parindent}{1.6em}
\setlist{nosep,leftmargin=2em}

% ---- 页眉 ----
\pagestyle{fancy}
\fancyhf{}
\fancyhead[L]{\small\coursename}
\fancyhead[R]{\small\thepage}
\renewcommand{\headrulewidth}{0.4pt}

% ---- 定理环境：中文用正体，标题加粗，按小节共用一条流水号 ----
\newtheoremstyle{cnthm}{0.7em}{0.7em}{\normalfont}{}{\bfseries}{}{0.5em}{}
\theoremstyle{cnthm}
\newtheorem{thm}{定理}[section]
\newtheorem{lem}[thm]{引理}
\newtheorem{prop}[thm]{命题}
\newtheorem{cor}[thm]{推论}
\newtheorem{defn}[thm]{定义}
\newtheorem{rem}[thm]{注}
\newtheorem{exa}[thm]{例}
\newtheorem{conj}[thm]{猜想}
\renewcommand{\proofname}{\bfseries 证明}
\providecommand{\coursename}{}
% ==STYLE-END==
\renewcommand{\coursename}{数学分析\,·\,1010}
\title{数学分析 1010}
\date{}
\begin{document}
\maketitle

\section{\ Recall}

\begin{thm}
单调上升且有上界的数列，极限一定存在。
\end{thm}

我们想证明
\[
\lim_{n\to+\infty}\left(1+\frac1n\right)^n
\]
存在，只需证明该数列单调递增且有上界。

\textbf{回顾：均值不等式}
\[
\sqrt[n]{a_1a_2\cdots a_n}\leq\frac{a_1+a_2+\cdots+a_n}{n},
\]
当且仅当 $a_1=\cdots=a_n$ 时取等。

\section{\ $e$}

\subsubsection*{单调递增}

令 $a_n=\left(1+\dfrac1n\right)^n$，要说明 $\{a_n\}$ 单增，只需证明
$\forall n,\ a_n\leq a_{n+1}$，即
\[
\left(1+\frac1n\right)^n\leq\left(1+\frac{1}{n+1}\right)^{n+1}.
\]
想用均值不等式，两边项数必须相同：给 LHS 乘上一个 $1$，
既不改变原式的值，又多出来了一项。
\[
\underbrace{\left(1+\frac1n\right)\cdots\left(1+\frac1n\right)}_{n\ \text{个}}\times1
\leq\left(\frac{n\left(1+\frac1n\right)+1}{n+1}\right)^{n+1}
=\left(1+\frac{1}{n+1}\right)^{n+1}.
\]
右边括号里的算术平均 $\dfrac{n\left(1+\frac1n\right)+1}{n+1}
=\dfrac{n+2}{n+1}=1+\dfrac{1}{n+1}$，正是要证的右边。
可知数列 $\{a_n\}$ 单调递增。

\subsubsection*{有界}

注意到 $b_n=\left(1+\dfrac1n\right)^{n+1}$ 单调递减，则对于所有 $n$，有
\[
a_n\leq b_n=a_n\left(1+\frac1n\right)\leq\cdots\leq b_1=4,
\]
可知整个数列是有界的。下补证 $\{b_n\}$ 单调递减：
\[
b_{n+1}\leq b_n
\iff\left(1+\frac{1}{n+1}\right)^{n+2}
\leq\left(1+\frac1n\right)^{n+1}.
\]
整理、取倒有
\[
\left(1-\frac{1}{n+1}\right)^{n+1}
\leq\left(1-\frac{1}{n+2}\right)^{n+2}.
\]
对 $n+2$ 个正数 $\underbrace{\left(1-\frac{1}{n+2}\right),\ldots,
\left(1-\frac{1}{n+2}\right)}_{n+1\ \text{个}},1$ 用均值不等式，得
\[
\left(1-\frac{1}{n+2}\right)^{n+1}\cdot1
\leq\left(\frac{(n+1)\left(1-\frac{1}{n+2}\right)+1}{n+2}\right)^{n+2}
=\left(\frac{n+1}{n+2}\right)^{n+2},
\]
两边开 $n+1$ 次方即得所需不等式，故 $\{b_n\}$ 单调递减。

由单调有界定理，$\displaystyle\lim\limits_{n\to+\infty}\left(1+\frac1n\right)^n$ 存在，
这个极限记作 $\mathrm e$。

\subsubsection*{例题：$1^\infty$ 型的极限}

下面通过几个例题来进一步阐释 $\mathrm e$。
\[
\lim_{n\to+\infty}\left(1+\frac{1}{n+1}\right)^n
=\lim_{n\to+\infty}\left(1+\frac{1}{n+1}\right)^{(n+1)\times\frac{n}{n+1}}
=\mathrm e .
\]
推而广之
\[
\lim_{n\to+\infty}\left(1+\frac{1}{n^2}\right)^{\frac{3n^3}{n+1}}
=\lim_{n\to+\infty}\left(1+\frac{1}{n^2}\right)^{n^2\times\frac{3n}{n+1}}
=\mathrm e^3 .
\]

\subsubsection*{Remark：数列趋向负无穷}

当数列趋向负无穷时，又会有什么性质呢？
\[
\lim_{n\to-\infty}\left(1+\frac1n\right)^n=?
\]
令 $m=-n$，有
\[
\begin{aligned}
\lim_{m\to+\infty}\left(1-\frac1m\right)^{-m}
&=\lim_{m\to+\infty}\left(\frac{m}{m-1}\right)^{-m}\\
&=\lim_{m\to+\infty}\left(1+\frac{1}{m-1}\right)^m\\
&=\lim_{m\to+\infty}\left(1+\frac{1}{m-1}\right)^{m-1}
\cdot\left(1+\frac{1}{m-1}\right)\\
&=\mathrm e .
\end{aligned}
\]
可见，极限也是 $\mathrm e$。

\section{\ 归一原理}

设 $\{a_n\}$ 极限存在，那对于 $\{a_{2k+1}\}$ 和 $\{a_{2k}\}$，
当 $k\to+\infty$ 时，是否存在极限？

反过来，若 $\{a_{2k+1}\}$ 和 $\{a_{2k}\}$，当 $k\to+\infty$ 时有相同的
极限，$a_n$ 是否存在极限？

\subsubsection*{正方向}

设 $\{a_n\}$ 有极限，即
\[
\lim_{n\to+\infty}a_n=A,
\]
则有
\[
\lim_{k\to+\infty}a_{2k}=A .
\]
由极限定义易得
$\forall\varepsilon>0$，$\exists N$，$\forall n>N$，
$|a_n-A|<\varepsilon$。我们要寻找 $K$，只要 $k>K$ 有
$|b_k-A|<\varepsilon$，只需 $2K>N$，那么
$\forall k>K$，
\[
|b_k-A|=|a_{2k}-A|<\varepsilon\qquad(2k>2K>N).
\]

\subsubsection*{反方向}

反过来，设
\[
\lim_{k\to+\infty}a_{2k}=\lim_{k\to+\infty}a_{2k+1}=A,
\]
下证明
\[
\lim_{n\to+\infty}a_n=A .
\]
由极限定义易得
\[
\begin{aligned}
\forall\varepsilon>0,\quad
&\exists K_1\ \text{s.t.}\ \forall k>K_1,\ |a_{2k}-A|<\varepsilon;\\
&\exists K_2\ \text{s.t.}\ \forall k>K_2,\ |a_{2k+1}-A|<\varepsilon .
\end{aligned}
\]
取 $K=\max\{K_1,K_2\}$，
\[
\begin{aligned}
\forall k>K,\ \text{有}\quad
&|a_{2k}-A|<\varepsilon,\\
&|a_{2k+1}-A|<\varepsilon .
\end{aligned}
\]
取 $N=2K$，则 $\forall n>N=2K$，无论 $n=2l$ or $2l+1$
都能得出 $l>K$，即得 $|a_n-A|<\varepsilon$。

\subsubsection*{推而广之}

数列 $\{a_n\}$ 有极限，等价于其任意子列都有极限。

设
\[
\lim_{n\to+\infty}a_n=A,
\]
根据极限定义翻译有
\[
\forall\varepsilon>0,\ \exists N_\varepsilon,\ \forall n>N_\varepsilon,\
\text{s.t.}\ |a_n-A|<\varepsilon .
\]
对于 $N_\varepsilon$ 总能找到 $k$ 使得 $n_k>N_\varepsilon$，那么有
\[
\forall k>K,\ n_k>N_\varepsilon,\ \text{s.t.}\ |a_{n_k}-A|<\varepsilon,
\]
即
\[
\lim_{k\to+\infty}a_{n_k}=A .
\]
反过来，特别地，数列 $\{a_n\}$ 本身就是自己的子列。

\end{document}

\end{document}
